\chapter{Florida Standards of Practice, FDOT and the New Datums}

This module ties mobile LiDAR to the rules a Florida Professional Surveyor and Mapper (PSM) works under: the \textbf{Standards of Practice in Chapter 5J-17, Florida Administrative Code (F.A.C.)}, the \textbf{FDOT Surveying and Mapping Handbook} requirements for Terrestrial Mobile LiDAR (TML), and the \textbf{modernization of the National Spatial Reference System (NSRS)}. Rule text is paraphrased from the versions cited in the sources box; always work from the current official text.

\begin{warning}[title={REGULATORY REVIEW REQUIRED BEFORE TEACHING}]
This module was compiled from the rule and agency texts listed at the end of the module (retrieved October 2026). It is \textbf{not legal advice} and has \textbf{not yet been reviewed by a second Florida PSM}. Before the seminar, a reviewing PSM should check every citation against the current official text and sign below.\\[4pt]
Reviewer: \_\_\_\_\_\_\_\_\_\_\_\_\_\_\_\_\_\_, PSM No.\ \_\_\_\_\_\_\_\_ \quad Date: \_\_\_\_\_\_\_\_ \quad Rule versions checked: \_\_\_\_\_\_\_\_\_\_\_\_
\end{warning}

\begin{kid}
Every game has a rulebook. Surveyors in Florida have one too, written by the state. It says what every map must show, how accurate it must be and how you prove it, and how you sign your name to it. When the map is made with a laser scanner instead of a total station, the rulebook still applies --- the scanner is just a new kind of measuring tape.
\end{kid}

\section{Who sets the rules}
\begin{center}\small
\renewcommand{\arraystretch}{1.25}
\begin{tabularx}{\linewidth}{>{\bfseries}p{3.6cm}X}\toprule
Chapter 472, Florida Statutes & the law that defines and licenses the practice of surveying and mapping\\
Board of Professional Surveyors and Mappers & adopts the rules; administratively housed in the Florida Department of Agriculture and Consumer Services\\
Chapter 5J-17, F.A.C. & the rules, including the Standards of Practice (5J-17.050 to 5J-17.053) and electronic signing and sealing (5J-17.062)\\
FDOT Surveying and Mapping Handbook & contract requirements for surveys on FDOT projects, including Terrestrial Mobile LiDAR (Sections 35--40)\\
NGS / Federal Register & the national datums (NSRS) that every Florida coordinate is tied to\\\bottomrule
\end{tabularx}\end{center}

\section{LiDAR is inside the Standards of Practice: definitions (5J-17.050)}
\begin{center}\small
\renewcommand{\arraystretch}{1.25}
\begin{tabularx}{\linewidth}{>{\bfseries}p{2.6cm}X}\toprule
5J-17.050(10) & \textbf{Remote sensing}: acquiring knowledge about physical features by analyzing data obtained from instrumentation --- a mobile LiDAR survey is remote sensing.\\
5J-17.050(11) & \textbf{Survey} includes control surveys and photogrammetric surveys, ``whether it is measured by direct or remote sensing methods.''\\
5J-17.050(7) & \textbf{Map of survey}: a graphical \emph{or digital} depiction --- a point cloud or digital surface delivered as a survey product is a map.\\
5J-17.050(8) & \textbf{Measurement records}: sufficient data to reconstruct the measurements --- for LiDAR, raw sensor data, trajectory and calibration.\\
5J-17.050(12) & \textbf{Survey and map report}: the map and report identify the responsible surveyor and mapper.\\
5J-17.050(13) & \textbf{Work product}: includes work made under the supervision of a PSM in responsible charge.\\\bottomrule
\end{tabularx}\end{center}

\section{General requirements applied to mobile LiDAR (5J-17.051)}
\begin{center}\small
\renewcommand{\arraystretch}{1.22}
\begin{tabularx}{\linewidth}{>{\bfseries}p{1.9cm}>{\raggedright\arraybackslash}p{4.6cm}X}\toprule
Rule & Requirement (paraphrased) & What it means for a mobile LiDAR deliverable\\\midrule
.051(2) & accuracy depends on survey type and intended use; \emph{the accuracy of data shall be independently verified} & test against \textbf{independent checkpoints} never used as control (Module 13); report RMSE in ASPRS Edition 2 terms\\
.051(2)(a) & products with elevation data state the \textbf{vertical datum} and describe the control; document vertical accuracy and the method used & e.g.\ NAVD 88 via GEOID18; list benchmarks/control; include the NVA test (30+ checkpoints)\\
.051(2)(b) & products showing mapped features state the \textbf{horizontal datum} (and control) & e.g.\ NAD 83(2011) epoch 2010.0, SPCS Florida East; FPRN mountpoint if RTK used\\
.051(2)(c) & if an assumed datum is used, adequate control points are established and communicated & avoid local-only SLAM frames; tie to control so the survey can be reproduced\\
.051(2)(d) & measurement and computation records dated to the acquisition date, with enough data to support the products & archive raw LiDAR/IMU/GNSS, trajectory, calibration file, processing logs (Section~\ref{sec:records})\\
.051(3)(a)--(b) & results communicated by a map, or a report with an attached map; type and/or purpose stated & e.g.\ ``Topographic Survey --- Mobile LiDAR'' with purpose (design, as-built)\\
.051(3)(c) & business name, certificate of authorization number, addresses; surveyor's name and license number & on the map, report and data index\\
.051(3)(d) & survey date = \textbf{last date of data acquisition}; revisions dated & the last collection pass date (not the processing date)\\
.051(3)(e) & unsigned copies state they are not valid without the signature and seal (electronic versions excepted) & see 5J-17.062 below\\
.051(3)(f)--(h) & abbreviations defined; bearing reference stated; north arrow and scale & applies to plotted LiDAR products and CAD sheets\\
.051(3)(i) & responsibility stated for all mapped features; if data from different surveyors are combined, identify who is primarily responsible & when LiDAR from another firm is merged with your field survey, say so\\\bottomrule
\end{tabularx}\end{center}

\begin{tip}
\textbf{5J-17.052 still governs boundaries.} A LiDAR topographic survey does not satisfy boundary-survey requirements (monuments, measured dimensions, water-boundary descriptions and so on in 5J-17.052). If a LiDAR product is part of a boundary survey, the boundary items must still be established and shown under 5J-17.052.
\end{tip}

\section{Signing and sealing digital LiDAR deliverables (5J-17.062)}
\begin{center}\small
\renewcommand{\arraystretch}{1.22}
\begin{tabularx}{\linewidth}{>{\bfseries}p{1.9cm}X}\toprule
.062(1) & electronic files representing documents required to be sealed under Ch.\ 472, F.S.\ are signed, dated and sealed by the PSM \textbf{in responsible charge}\\
.062(2) & a computer-generated seal may be used; final hard copies carry an original signature, seal and date (or an electronic signature); a \textbf{scanned signature is not acceptable}\\
.062(3) & an \textbf{electronic signature} has the same effect as an original signature and seal if it is unique to the user, verifiable, under the user's sole control, and \textbf{invalidated if any data in the document change}\\
.062(4) & alternatively, a \textbf{signature file} lists the PSM's name and number, a description, and each file's name and \textbf{secure authentication code} (a cryptographic hash); a printed report of it is manually signed, dated and sealed\\\bottomrule
\end{tabularx}\end{center}

\begin{worked}{sealing a 40-tile LAZ delivery with a signature file}
\step{Step 1} Compute a cryptographic hash (e.g.\ SHA-256) for every delivered file: 40 LAZ tiles, the TIN, the CAD file, the trajectory file and the PDF report.
\step{Step 2} Build the signature file: PSM name and number, project description, and one line per file --- \texttt{tile\_0101.laz  SHA-256: 9f2c\ldots e41a} (64 hexadecimal characters each).
\step{Step 3} Print the signature file report; sign, date and seal it manually.
\step{Step 4 --- verification} anyone can recompute each file's hash; if it matches the sealed list, that file is the sealed file. If one point in one tile is edited, that tile's hash changes completely and it is no longer sealed.
\result{this is how a multi-gigabyte point cloud can be sealed without printing it --- and why every later edit must be re-delivered with a new sealed signature file.}
\end{worked}

\section{FDOT Terrestrial Mobile LiDAR (TML) requirements}
\begin{eng}
On FDOT work, TML surveys follow the \textbf{FDOT Surveying and Mapping Handbook} (current edition dated February 28, 2025, superseding the 3/29/2019 edition). Its TML part is organized as follows; read these sections before scoping a project:
\begin{center}\small
\renewcommand{\arraystretch}{1.15}
\begin{tabularx}{\linewidth}{p{3.2cm}X}\toprule
\textbf{Section} & \textbf{Topic}\\\midrule
35 & Introduction\\
36, 36.1 & TML methods; types of TML surveys\\
37 & TML project selection\\
38, 38.1 & TML equipment; minimum TML system sensor components\\
39.1--39.5 & mission; base station control; maintenance and boresight calibration; redundancy; monitoring\\
39.6, 39.6.1, 39.6.2 & project control and checkpoints; Type ``A'' control layout; TML survey specifications\\
39.7, 39.8 & accuracy analysis; quality assurance/quality control\\
40, 40.1, 40.1.1, 40.2 & deliverables; all-type deliverables; LiDAR point classes; TML survey report\\
Appendix C & accuracy requirements\\\bottomrule
\end{tabularx}\end{center}
FDOT defines three survey types by accuracy (originally in its 2012 TML guidelines, adapted from Caltrans): \textbf{Type A} (high accuracy) for design engineering and as-built surveys; \textbf{Type B} (medium) for detailed asset inventory and earthwork; \textbf{Type C} (low) for planning and transportation statistics.
\end{eng}

\begin{center}\small
\renewcommand{\arraystretch}{1.3}
\begin{tabularx}{\linewidth}{>{\bfseries}p{1.6cm}Xp{3.2cm}p{3.2cm}}\toprule
Type & Typical use & Accuracy requirement (enter from Handbook §39.6.2 / App.\ C) & Control/checkpoint layout (enter from §39.6, §39.6.1)\\\midrule
A & design engineering, as-built & \_\_\_\_\_\_\_\_ & \_\_\_\_\_\_\_\_\\
B & asset inventory, earthwork & \_\_\_\_\_\_\_\_ & \_\_\_\_\_\_\_\_\\
C & planning, statistics & \_\_\_\_\_\_\_\_ & \_\_\_\_\_\_\_\_\\\bottomrule
\end{tabularx}\end{center}
{\footnotesize The numeric tolerances are deliberately left blank: copy them from the Handbook edition named in your contract so the seminar never teaches an outdated number.}

\begin{tip}
\textbf{District practice example (FDOT District 5 QA/QC plan).} ``Imagery'' covers aerial photogrammetry \emph{and} TML. Coordinate with the FDOT Central Office for a Mobile Survey Tracking System (MSTS) number; submit all raw, interim and final imagery data to the Central Office with a letter of transmittal; deliver the derived DGN/TIN files to the District; and note in the Survey and Map Report that the Central Office has received a separate certified report covering the imagery data. Check your own District's plan.
\end{tip}

\section{The new national datums (NSRS modernization)}\label{sec:nsrs}
\fig[0.95]{nsrs_timeline}{NSRS modernization timeline as published by NGS and the Federal Register. Dates after 2026 are NGS expectations, not fixed.}

\begin{eng}
NGS is replacing NAD 83 and NAVD 88 with a modernized NSRS:
\begin{itemize}
\item \textbf{NATRF2022} (North American Terrestrial Reference Frame of 2022), with PATRF2022, CATRF2022 and MATRF2022 for other plates --- replacing the NAD 83 realizations.
\item \textbf{NAPGD2022} (North American-Pacific Geopotential Datum of 2022), including \textbf{GEOID2022} --- replacing NAVD 88 and the hybrid geoids such as GEOID18.
\item \textbf{SPCS2022} --- new State Plane zones, defined in meters with the \textbf{international foot} as the only foot.
\end{itemize}
\textbf{Timeline.} The Federal Register notice of October 9, 2024 scheduled the modernization for 2025 or 2026, with components posted for public beta testing (at least six months after the last component), a vote by the Federal Geodetic Control Subcommittee (FGCS), FGDC endorsement and a Federal Register adoption notice, followed by tools and services within five years. NGS's February 2026 update placed the FGCS vote in \textbf{late 2026 or early 2027} and said full adoption \textbf{may extend into the first part of 2027}. \textbf{Until adoption, the current NSRS (NAD 83, NAVD 88) remains official.}

\textbf{Units.} NIST and NOAA deprecated the \textbf{U.S. survey foot on December 31, 2022} in favor of the international foot (0.3048 m exactly). The difference is 2 parts per million --- about 0.2 ft over the 100\,000 ft magnitudes of Florida State Plane coordinates. Section 177.151, F.S.\ expresses Florida State Plane coordinates ``in feet and decimals of a foot'': state explicitly which foot each deliverable uses, and confirm the current statute and the client's requirement.
\end{eng}

\begin{worked}{why the foot matters}
\textbf{Given:} a State Plane easting of 650\,000 ft.
\step{U.S. survey foot} $1\text{ ft}=1200/3937$ m $=0.3048006096$ m $\Rightarrow$ $650\,000\times0.3048006096=198\,120.396$ m
\step{International foot} $650\,000\times0.3048=198\,120.000$ m
\step{Difference} 0.396 m $=$ 1.30 ft
\result{a mislabeled foot moves every point by more than a foot --- far larger than any LiDAR error in this course. Always label the foot.}
\end{worked}

\begin{tip}
\textbf{Future-proof today's LiDAR.} Record datum, realization, epoch, geoid model and foot in the report, the LAS header (WKT) and the metadata; keep ellipsoid heights and the trajectory so the project can be transformed to NATRF2022/NAPGD2022 later; and expect clients to ask for re-delivery in the new datums once they are adopted.
\end{tip}

\section{Checklist: a Florida-compliant mobile LiDAR deliverable}
\begin{center}\small
\renewcommand{\arraystretch}{1.18}
\begin{tabularx}{\linewidth}{lX}\toprule
$\square$ & map or report with attached map; type and purpose stated (.051(3)(a)--(b))\\
$\square$ & business, certificate of authorization, PSM name and license (.051(3)(c))\\
$\square$ & survey date = last acquisition date; revisions dated (.051(3)(d))\\
$\square$ & horizontal datum, realization, epoch, SPCS zone, foot; vertical datum and geoid; control described (.051(2)(a)--(b))\\
$\square$ & accuracy independently verified with checkpoints; method documented; ASPRS statement (.051(2))\\
$\square$ & measurement and computation records archived and dated (.051(2)(d))\\
$\square$ & responsibility for merged data stated (.051(3)(i))\\
$\square$ & electronic files sealed by electronic signature or sealed signature file (.062)\\
$\square$ & FDOT projects: Handbook TML type, control/checkpoint layout, report and Central Office submittal\\\bottomrule
\end{tabularx}\end{center}

\begin{tryit}
\textbf{19.1} (Concept) Which 5J-17.051 subsection requires that the accuracy of data be independently verified, and how does a LiDAR survey satisfy it?\\
\textbf{19.2} (Concept) What is the ``survey date'' of a mobile LiDAR project collected on March 3 and March 5 and processed on March 20?\\
\textbf{19.3} (Concept) A client edits one LAZ tile after delivery. Is it still sealed under a 5J-17.062(4) signature file? Why?\\
\textbf{19.4} Convert a 1\,200\,000 ft State Plane northing using both feet. How large is the difference?
\end{tryit}

\begin{tcolorbox}[base,colback=gray!4,colframe=grayx,title={SOURCES FOR THIS MODULE}]
\footnotesize
Fla.\ Admin.\ Code R.\ 5J-17.050, 5J-17.051 (eff.\ 11/13/2017), 5J-17.052 and 5J-17.062 (eff.\ 7/15/2018), as published by Cornell LII; Ch.\ 472 and s.\ 177.151, F.S.; FDOT \emph{Surveying and Mapping Handbook} (Feb.\ 28, 2025), Sections 35--40 and App.\ C; FDOT \emph{Terrestrial Mobile LiDAR Surveying \& Mapping Guidelines} (2012) as summarized by LiDAR News; FDOT District 5 \emph{Survey QA/QC Management Plan}; Federal Register 2024-23347, ``Updated Implementation Timeline for the Modernized NSRS'' (Oct.\ 9, 2024); GPS World, ``NGS presents the latest NSRS news at Geo Week 2026'' (Mar.\ 3, 2026); Esri ArcNews, ``Change is afoot with the US survey foot.''
\end{tcolorbox}
